\chapter{Cardinality, Normal Forms, and Resolution}
\label{ch:fol-cardinality-normalforms}

\section{Elementary substructures}
\label{sec:fol-cardinality-normalforms-01}

\begin{definition}[Substructure]
Let \(\Meta\) and \(\mathcal N\) be \(\Lang\)-structures.  Then
\(\mathcal N\) is a \emph{substructure} of \(\Meta\), written
\(\mathcal N\subseteq\Meta\), iff:
\begin{enumerate}
  \item \(\dom\mathcal N\subseteq\dom\Meta\);
  \item every constant has the same value in both structures;
  \item \(\dom\mathcal N\) is closed under every function of \(\Meta\),
        and each function of \(\mathcal N\) is the corresponding restriction;
        and
  \item every relation of \(\mathcal N\) is the restriction of the
        corresponding relation of \(\Meta\) to \(\dom\mathcal N\).
\end{enumerate}
\end{definition}

\begin{definition}[Elementary substructure]
A substructure \(\mathcal N\subseteq\Meta\) is an \emph{elementary
substructure}, written \(\mathcal N\preccurlyeq\Meta\), iff for every formula
\(\varphi(x_1,\ldots,x_n)\) and every
\(a_1,\ldots,a_n\in\dom\mathcal N\),
\[
 \mathcal N\entails\varphi(a_1,\ldots,a_n)
 \quad\text{iff}\quad
 \Meta\entails\varphi(a_1,\ldots,a_n).
\]
The notation with parameters abbreviates satisfaction under assignments that
map the displayed variables to the displayed elements.
\end{definition}

\begin{theorem}[Tarski--Vaught test]
\label{thm:tarski-vaught-fol}
Let \(\mathcal N\subseteq\Meta\).  Then
\(\mathcal N\preccurlyeq\Meta\) iff, for every formula
\(\varphi(y,x_1,\ldots,x_n)\) and every tuple
\(\bar a\in(\dom\mathcal N)^{n}\),
\[
 \Meta\entails\exists y\,\varphi(y,\bar a)
 \quad\Longrightarrow\quad
 \Meta\entails\varphi(b,\bar a)
 \text{ for some }b\in\dom\mathcal N.
\]
\end{theorem}
\begin{proof}
If \(\mathcal N\preccurlyeq\Meta\) and the existential formula with
parameters is satisfied in \(\Meta\), it is satisfied in
\(\mathcal N\); its witness
belongs to \(\dom\mathcal N\).

Conversely, assume the witness condition and prove elementarity by structural
induction on formulas.  Term values of tuples from \(\mathcal N\) agree in
the two structures because \(\mathcal N\) is closed under the functions.
Atomic formulas therefore agree because relations are restricted.  The
Boolean cases follow from the induction hypotheses.  If an existential formula
with parameters from \(\mathcal N\) is true in \(\mathcal N\), its witness is
also in \(\Meta\), and the induction hypothesis transfers the matrix.  If it
is true in \(\Meta\), the assumed condition supplies a witness in
\(\mathcal N\), and the induction hypothesis transfers the matrix back.
For a universal formula \(\forall y\,\varphi\) with parameters from
\(\mathcal N\): if it is true in \(\Meta\), then every instance at an
element of \(\dom\mathcal N\) is true in \(\Meta\), and the induction
hypothesis for \(\varphi\) transfers each instance to \(\mathcal N\).  If
it is false in \(\Meta\), the satisfaction clauses give
\(\Meta\entails\exists y\,\neg\varphi\) at the same parameters; the
witness condition supplies a witness in \(\dom\mathcal N\), the
induction hypothesis for \(\varphi\) transfers the negated instance to
\(\mathcal N\), and \(\forall y\,\varphi\) is false in
\(\mathcal N\).
\end{proof}

